% Sección del informe: evolución del extractor de etiquetas débiles (entregables 1 y 7).
% Uso en el main.tex de NeurIPS:  % Sección del informe: evolución del extractor de etiquetas débiles (entregables 1 y 7).
% Uso en el main.tex de NeurIPS:  % Sección del informe: evolución del extractor de etiquetas débiles (entregables 1 y 7).
% Uso en el main.tex de NeurIPS:  % Sección del informe: evolución del extractor de etiquetas débiles (entregables 1 y 7).
% Uso en el main.tex de NeurIPS:  \input{evolucion_extractor}   (requiere \usepackage{booktabs}, ya incluido en la plantilla)
% Fuente de los números: salida de MODE='labels' en Kaggle (4407 reportes, 58 estudios gold × 12 hallazgos = 696 decisiones).

\subsection{Extractor de etiquetas débiles: validación y evolución}
\label{sec:extractor}

Solo 58 de los 4407 estudios de entrenamiento tienen etiqueta experta (\emph{gold}).
El resto se etiqueta con un extractor propio de reglas que lee el reporte radiológico, escrito en al menos nueve idiomas.
Para cada hallazgo, el extractor busca menciones (término directo, o anatomía + patología a menos de 60 caracteres) y clasifica cada una según su contexto inmediato, al estilo de NegEx~\cite{chapman2001negex}:
negación $\rightarrow 0$, duda $\rightarrow 0.5$, edema óseo sin la palabra \emph{contusión} $\rightarrow 0.4$, condropatía de bajo grado $\rightarrow 0.4$, derrame leve $\rightarrow 0.5$, y afirmación $\rightarrow 1$.
El puntaje del estudio es el máximo de sus menciones.
Si no hay ninguna mención (\emph{silencio}), la etiqueta queda indefinida y en el entrenamiento recibe peso $0.3$.

\paragraph{Métricas.}
Cada versión se compara con los 58 gold.
El \emph{acuerdo} es la fracción de decisiones binarias (puntaje $\geq 0.5$, silencio $=0$) que coinciden con el experto, y la \emph{tasa de ruido} es $1-\text{acuerdo}$.
La \emph{mención} es la fracción de reportes gold en que el extractor encontró algo sobre el hallazgo.
El AUC usa el puntaje continuo y no depende del umbral, por eso es la métrica principal para comparar versiones.

\paragraph{Evolución.}
El extractor pasó por cuatro versiones (Tabla~\ref{tab:ext-cambios}).
Cada cambio salió de leer reportes reales con la función \texttt{show\_reports}, no de ajustar casos gold puntuales.
Cada regla nueva quedó fijada con una prueba de regresión: 27 en total.
El AUC macro subió de $0.739$ a $0.753$ y la tasa de ruido bajó de $0.226$ a $0.218$ (Tabla~\ref{tab:ext-versiones}).
La mención media pasó de $0.509$ a $0.566$ (Tabla~\ref{tab:ext-final}).

\begin{table}[t]
  \caption{Cambios del extractor por versión. Todos los cambios provienen de leer reportes reales.}
  \label{tab:ext-cambios}
  \centering
  \small
  \begin{tabular}{@{}lp{0.42\linewidth}p{0.42\linewidth}@{}}
    \toprule
    Versión & Cambio & Motivación \\
    \midrule
    v1 & Léxico multilingüe anatomía + patología; negación y duda estilo NegEx & Línea base \\
    v2 & Derrame y sinovitis leves $\rightarrow 0.3$; edema óseo $\rightarrow 0.4$; sinónimos de artrosis (cartílago, osteofitos, pinzamiento); sinovitis por engrosamiento sinovial y grasa de Hoffa; rotura meniscal por ``señal que alcanza la superficie''; puntaje del estudio = máximo
      & Especificidad de derrame $0.35$; sensibilidad de OA lateral y patelofemoral $<0.40$; mención de sinovitis $0.26$ \\
    v3 & Derrame leve $\rightarrow 0.5$ (incierto); sinovitis leve $\rightarrow 1$; condropatía grado I--II o superficial $\rightarrow 0.4$
      & En v2, los derrames leves resultaron positivos en el gold en proporción cercana a la mitad; grado I--II de Outerbridge no es artrosis \\
    v4 & Negación de estructuras nombradas en conjunto (``meniscos'', ``ligamentos colaterales'', ``cartílago'', ``médula ósea''); unión de líneas cortadas; ventana de negación de 35 a 60 caracteres; ``grado 2 y 3'' cuenta como alto grado; un lado anatómico solo (``menisco lateral íntegro'') ya no niega la artrosis de ese lado; ``osteocondral'' en griego deja de contar como artrosis
      & Reportes en español, griego y neerlandés que listan lo normal en bloque quedaban en silencio; falsos negativos y falsos positivos de artrosis \\
    \bottomrule
  \end{tabular}
\end{table}

\begin{table}[t]
  \caption{Acuerdo y AUC del extractor frente a los 58 gold, por versión. $n^{+}$: positivos según el experto.}
  \label{tab:ext-versiones}
  \centering
  \small
  \begin{tabular}{@{}lrcccccccc@{}}
    \toprule
    & & \multicolumn{4}{c}{Acuerdo} & \multicolumn{4}{c}{AUC} \\
    \cmidrule(lr){3-6} \cmidrule(l){7-10}
    Hallazgo & $n^{+}$ & v1 & v2 & v3 & v4 & v1 & v2 & v3 & v4 \\
    \midrule
    ACL              & 24 & .879 & .879 & .879 & .879 & .885 & .885 & .885 & .885 \\
    MCL              &  9 & .828 & .828 & .828 & .828 & .754 & .754 & .754 & .754 \\
    Medial Meniscus  & 26 & .776 & .776 & .776 & .759 & .768 & .768 & .768 & .752 \\
    Lateral Meniscus & 23 & .828 & .828 & .828 & .828 & .806 & .806 & .806 & .806 \\
    Medial OA        & 15 & .828 & .776 & .776 & .810 & .747 & .740 & .740 & .760 \\
    Lateral OA       & 11 & .810 & .793 & .810 & .810 & .635 & .725 & .733 & .733 \\
    PF OA            & 21 & .707 & .672 & .672 & .707 & .636 & .655 & .655 & .690 \\
    Effusion         & 35 & .638 & .621 & .638 & .638 & .666 & .666 & .666 & .666 \\
    Synovitis        & 27 & .655 & .621 & .672 & .672 & .636 & .655 & .658 & .658 \\
    Baker's          & 12 & .862 & .862 & .862 & .862 & .821 & .821 & .821 & .821 \\
    Contusion        & 19 & .690 & .793 & .793 & .793 & .751 & .760 & .760 & .760 \\
    Fracture         & 18 & .793 & .793 & .793 & .793 & .758 & .758 & .758 & .758 \\
    \midrule
    Media / macro    &    & .774 & .770 & .777 & .782 & .739 & .749 & .750 & .753 \\
    Tasa de ruido    &    & .226 & .230 & .223 & .218 &      &      &      &      \\
    \bottomrule
  \end{tabular}
\end{table}

\begin{table}[t]
  \caption{Versión final (v4) frente a los gold, con la mención de v1 como referencia.
    Sensibilidad y especificidad usan el umbral $0.5$.}
  \label{tab:ext-final}
  \centering
  \small
  \begin{tabular}{@{}lrcccccc@{}}
    \toprule
    & & \multicolumn{2}{c}{Mención} & & & & \\
    \cmidrule(lr){3-4}
    Hallazgo & $n^{+}$ & v1 & v4 & Acuerdo & Sens. & Espec. & AUC \\
    \midrule
    ACL              & 24 & .690 & .724 & .879 & .917 & .853 & .885 \\
    MCL              &  9 & .500 & .621 & .828 & .667 & .857 & .754 \\
    Medial Meniscus  & 26 & .517 & .621 & .759 & .692 & .812 & .752 \\
    Lateral Meniscus & 23 & .569 & .586 & .828 & .652 & .943 & .806 \\
    Medial OA        & 15 & .431 & .500 & .810 & .667 & .860 & .760 \\
    Lateral OA       & 11 & .448 & .448 & .810 & .636 & .851 & .733 \\
    PF OA            & 21 & .552 & .621 & .707 & .667 & .730 & .690 \\
    Effusion         & 35 & .845 & .845 & .638 & .829 & .348 & .666 \\
    Synovitis        & 27 & .259 & .362 & .672 & .444 & .871 & .658 \\
    Baker's          & 12 & .362 & .362 & .862 & .750 & .891 & .821 \\
    Contusion        & 19 & .517 & .603 & .793 & .684 & .846 & .760 \\
    Fracture         & 18 & .414 & .500 & .793 & .667 & .850 & .758 \\
    \midrule
    Media            &    & .509 & .566 & .782 & .689 & .809 & .753 \\
    \bottomrule
  \end{tabular}
\end{table}

\paragraph{Lectura por hallazgo.}
Los hallazgos que el reporte nombra de forma explícita y binaria son los mejor recuperados: ACL (AUC $0.885$), quiste de Baker ($0.821$) y menisco lateral ($0.806$).
El derrame muestra el límite de un umbral fijo.
De v1 a v3 su AUC no cambió ($0.666$), pero sensibilidad y especificidad se intercambiaron según el valor asignado al derrame \emph{leve}.
En v2 ese valor fue $0.3$: se perdieron 11 positivos del experto y se ganaron 10 negativos, es decir, los expertos consideran positiva aproximadamente la mitad de los derrames leves.
Por eso, desde v3 el derrame leve vale $0.5$ (incierto) y entra al entrenamiento con peso $0.5$.
La sinovitis tiene la mención más baja ($0.36$) aunque el experto la marca en 27 de 58 estudios.
Probablemente el experto la infiere de la imagen y no del texto, y ninguna regla sobre el reporte puede recuperar esa información.
v4 introdujo una regresión pequeña en el menisco medial: un estudio negativo pasó a positivo (especificidad de $0.844$ a $0.812$).

\paragraph{Sesgo lingüístico.}
La Tabla~\ref{tab:ext-idioma} desglosa la calidad por idioma.
El \emph{silencio} es la fracción de los 12 hallazgos sin mención y no requiere etiquetas.
En v1, un reporte en español quedaba en silencio en el $73\%$ de los hallazgos, frente al $45\%$ en inglés.
La diferencia venía sobre todo del estilo de redacción, no del vocabulario: los reportes en español, griego y neerlandés suelen declarar lo normal en bloque (``ligamentos cruzados y colaterales dentro de límites normales'', ``meniscos conservados'').
v4 redujo el silencio en español de $0.717$ a $0.560$.
En turco el silencio subió ($0.617 \rightarrow 0.703$) porque se eliminaron negativos espurios de artrosis.
En v3, palabras de lado como \emph{iç}/\emph{dış} (interno/externo) junto a ``normal'' negaban la artrosis de ese compartimento aunque la frase hablara del menisco.
Ese aumento de silencio es, por tanto, una corrección.
El acuerdo por idioma en griego ($0.53$) y neerlandés ($0.63$) se apoya en solo 2--3 estudios gold, unas 30 decisiones, y no es estadísticamente concluyente.
Además, algunos reportes neerlandeses tienen la sección de hallazgos vacía; ese silencio proviene del dato y no del extractor.

\begin{table}[t]
  \caption{Silencio medio (fracción de hallazgos sin mención) y acuerdo con los gold, por idioma del reporte.}
  \label{tab:ext-idioma}
  \centering
  \small
  \begin{tabular}{@{}lrrcccccc@{}}
    \toprule
    & & & \multicolumn{4}{c}{Silencio} & \multicolumn{2}{c}{Acuerdo gold} \\
    \cmidrule(lr){4-7} \cmidrule(l){8-9}
    Idioma & Reportes & Gold & v1 & v2 & v3 & v4 & v1 & v4 \\
    \midrule
    Inglés     & 1736 & 28 & .445 & .393 & .393 & .383 & .827 & .812 \\
    Español    &  682 & 10 & .730 & .717 & .717 & .560 & .783 & .842 \\
    Turco      &  546 &  6 & .633 & .617 & .617 & .703 & .750 & .764 \\
    Croata     &  407 &  4 & .598 & .527 & .527 & .509 & .729 & .729 \\
    Griego     &  321 &  3 & .781 & .737 & .737 & .661 & .528 & .528 \\
    Alemán     &  262 &  2 & .589 & .543 & .543 & .513 & .792 & .792 \\
    Cirílico   &  220 &  3 & .717 & .678 & .678 & .649 & .750 & .750 \\
    Neerlandés &  151 &  2 & .726 & .670 & .670 & .604 & .542 & .625 \\
    Francés    &   81 &  0 & .384 & .340 & .340 & .294 & --   & --   \\
    \bottomrule
  \end{tabular}
\end{table}

\paragraph{Limitaciones.}
Los 58 gold cumplen dos funciones: miden la calidad del extractor y, más adelante, evalúan el modelo final.
Dos parámetros del extractor (derrame leve $=0.5$ y sinovitis leve $=1$) se fijaron mirando su efecto sobre esos gold, lo que puede hacer algo optimista la evaluación final.
Para acotar ese sesgo, el resto de las reglas son clínicas y generales: grados de Outerbridge, negación de estructuras en bloque y sinónimos.
Además, la selección de época y de umbrales del modelo se hace en un \emph{holdout} con etiquetas débiles, no en los gold.
   (requiere \usepackage{booktabs}, ya incluido en la plantilla)
% Fuente de los números: salida de MODE='labels' en Kaggle (4407 reportes, 58 estudios gold × 12 hallazgos = 696 decisiones).

\subsection{Extractor de etiquetas débiles: validación y evolución}
\label{sec:extractor}

Solo 58 de los 4407 estudios de entrenamiento tienen etiqueta experta (\emph{gold}).
El resto se etiqueta con un extractor propio de reglas que lee el reporte radiológico, escrito en al menos nueve idiomas.
Para cada hallazgo, el extractor busca menciones (término directo, o anatomía + patología a menos de 60 caracteres) y clasifica cada una según su contexto inmediato, al estilo de NegEx~\cite{chapman2001negex}:
negación $\rightarrow 0$, duda $\rightarrow 0.5$, edema óseo sin la palabra \emph{contusión} $\rightarrow 0.4$, condropatía de bajo grado $\rightarrow 0.4$, derrame leve $\rightarrow 0.5$, y afirmación $\rightarrow 1$.
El puntaje del estudio es el máximo de sus menciones.
Si no hay ninguna mención (\emph{silencio}), la etiqueta queda indefinida y en el entrenamiento recibe peso $0.3$.

\paragraph{Métricas.}
Cada versión se compara con los 58 gold.
El \emph{acuerdo} es la fracción de decisiones binarias (puntaje $\geq 0.5$, silencio $=0$) que coinciden con el experto, y la \emph{tasa de ruido} es $1-\text{acuerdo}$.
La \emph{mención} es la fracción de reportes gold en que el extractor encontró algo sobre el hallazgo.
El AUC usa el puntaje continuo y no depende del umbral, por eso es la métrica principal para comparar versiones.

\paragraph{Evolución.}
El extractor pasó por cuatro versiones (Tabla~\ref{tab:ext-cambios}).
Cada cambio salió de leer reportes reales con la función \texttt{show\_reports}, no de ajustar casos gold puntuales.
Cada regla nueva quedó fijada con una prueba de regresión: 27 en total.
El AUC macro subió de $0.739$ a $0.753$ y la tasa de ruido bajó de $0.226$ a $0.218$ (Tabla~\ref{tab:ext-versiones}).
La mención media pasó de $0.509$ a $0.566$ (Tabla~\ref{tab:ext-final}).

\begin{table}[t]
  \caption{Cambios del extractor por versión. Todos los cambios provienen de leer reportes reales.}
  \label{tab:ext-cambios}
  \centering
  \small
  \begin{tabular}{@{}lp{0.42\linewidth}p{0.42\linewidth}@{}}
    \toprule
    Versión & Cambio & Motivación \\
    \midrule
    v1 & Léxico multilingüe anatomía + patología; negación y duda estilo NegEx & Línea base \\
    v2 & Derrame y sinovitis leves $\rightarrow 0.3$; edema óseo $\rightarrow 0.4$; sinónimos de artrosis (cartílago, osteofitos, pinzamiento); sinovitis por engrosamiento sinovial y grasa de Hoffa; rotura meniscal por ``señal que alcanza la superficie''; puntaje del estudio = máximo
      & Especificidad de derrame $0.35$; sensibilidad de OA lateral y patelofemoral $<0.40$; mención de sinovitis $0.26$ \\
    v3 & Derrame leve $\rightarrow 0.5$ (incierto); sinovitis leve $\rightarrow 1$; condropatía grado I--II o superficial $\rightarrow 0.4$
      & En v2, los derrames leves resultaron positivos en el gold en proporción cercana a la mitad; grado I--II de Outerbridge no es artrosis \\
    v4 & Negación de estructuras nombradas en conjunto (``meniscos'', ``ligamentos colaterales'', ``cartílago'', ``médula ósea''); unión de líneas cortadas; ventana de negación de 35 a 60 caracteres; ``grado 2 y 3'' cuenta como alto grado; un lado anatómico solo (``menisco lateral íntegro'') ya no niega la artrosis de ese lado; ``osteocondral'' en griego deja de contar como artrosis
      & Reportes en español, griego y neerlandés que listan lo normal en bloque quedaban en silencio; falsos negativos y falsos positivos de artrosis \\
    \bottomrule
  \end{tabular}
\end{table}

\begin{table}[t]
  \caption{Acuerdo y AUC del extractor frente a los 58 gold, por versión. $n^{+}$: positivos según el experto.}
  \label{tab:ext-versiones}
  \centering
  \small
  \begin{tabular}{@{}lrcccccccc@{}}
    \toprule
    & & \multicolumn{4}{c}{Acuerdo} & \multicolumn{4}{c}{AUC} \\
    \cmidrule(lr){3-6} \cmidrule(l){7-10}
    Hallazgo & $n^{+}$ & v1 & v2 & v3 & v4 & v1 & v2 & v3 & v4 \\
    \midrule
    ACL              & 24 & .879 & .879 & .879 & .879 & .885 & .885 & .885 & .885 \\
    MCL              &  9 & .828 & .828 & .828 & .828 & .754 & .754 & .754 & .754 \\
    Medial Meniscus  & 26 & .776 & .776 & .776 & .759 & .768 & .768 & .768 & .752 \\
    Lateral Meniscus & 23 & .828 & .828 & .828 & .828 & .806 & .806 & .806 & .806 \\
    Medial OA        & 15 & .828 & .776 & .776 & .810 & .747 & .740 & .740 & .760 \\
    Lateral OA       & 11 & .810 & .793 & .810 & .810 & .635 & .725 & .733 & .733 \\
    PF OA            & 21 & .707 & .672 & .672 & .707 & .636 & .655 & .655 & .690 \\
    Effusion         & 35 & .638 & .621 & .638 & .638 & .666 & .666 & .666 & .666 \\
    Synovitis        & 27 & .655 & .621 & .672 & .672 & .636 & .655 & .658 & .658 \\
    Baker's          & 12 & .862 & .862 & .862 & .862 & .821 & .821 & .821 & .821 \\
    Contusion        & 19 & .690 & .793 & .793 & .793 & .751 & .760 & .760 & .760 \\
    Fracture         & 18 & .793 & .793 & .793 & .793 & .758 & .758 & .758 & .758 \\
    \midrule
    Media / macro    &    & .774 & .770 & .777 & .782 & .739 & .749 & .750 & .753 \\
    Tasa de ruido    &    & .226 & .230 & .223 & .218 &      &      &      &      \\
    \bottomrule
  \end{tabular}
\end{table}

\begin{table}[t]
  \caption{Versión final (v4) frente a los gold, con la mención de v1 como referencia.
    Sensibilidad y especificidad usan el umbral $0.5$.}
  \label{tab:ext-final}
  \centering
  \small
  \begin{tabular}{@{}lrcccccc@{}}
    \toprule
    & & \multicolumn{2}{c}{Mención} & & & & \\
    \cmidrule(lr){3-4}
    Hallazgo & $n^{+}$ & v1 & v4 & Acuerdo & Sens. & Espec. & AUC \\
    \midrule
    ACL              & 24 & .690 & .724 & .879 & .917 & .853 & .885 \\
    MCL              &  9 & .500 & .621 & .828 & .667 & .857 & .754 \\
    Medial Meniscus  & 26 & .517 & .621 & .759 & .692 & .812 & .752 \\
    Lateral Meniscus & 23 & .569 & .586 & .828 & .652 & .943 & .806 \\
    Medial OA        & 15 & .431 & .500 & .810 & .667 & .860 & .760 \\
    Lateral OA       & 11 & .448 & .448 & .810 & .636 & .851 & .733 \\
    PF OA            & 21 & .552 & .621 & .707 & .667 & .730 & .690 \\
    Effusion         & 35 & .845 & .845 & .638 & .829 & .348 & .666 \\
    Synovitis        & 27 & .259 & .362 & .672 & .444 & .871 & .658 \\
    Baker's          & 12 & .362 & .362 & .862 & .750 & .891 & .821 \\
    Contusion        & 19 & .517 & .603 & .793 & .684 & .846 & .760 \\
    Fracture         & 18 & .414 & .500 & .793 & .667 & .850 & .758 \\
    \midrule
    Media            &    & .509 & .566 & .782 & .689 & .809 & .753 \\
    \bottomrule
  \end{tabular}
\end{table}

\paragraph{Lectura por hallazgo.}
Los hallazgos que el reporte nombra de forma explícita y binaria son los mejor recuperados: ACL (AUC $0.885$), quiste de Baker ($0.821$) y menisco lateral ($0.806$).
El derrame muestra el límite de un umbral fijo.
De v1 a v3 su AUC no cambió ($0.666$), pero sensibilidad y especificidad se intercambiaron según el valor asignado al derrame \emph{leve}.
En v2 ese valor fue $0.3$: se perdieron 11 positivos del experto y se ganaron 10 negativos, es decir, los expertos consideran positiva aproximadamente la mitad de los derrames leves.
Por eso, desde v3 el derrame leve vale $0.5$ (incierto) y entra al entrenamiento con peso $0.5$.
La sinovitis tiene la mención más baja ($0.36$) aunque el experto la marca en 27 de 58 estudios.
Probablemente el experto la infiere de la imagen y no del texto, y ninguna regla sobre el reporte puede recuperar esa información.
v4 introdujo una regresión pequeña en el menisco medial: un estudio negativo pasó a positivo (especificidad de $0.844$ a $0.812$).

\paragraph{Sesgo lingüístico.}
La Tabla~\ref{tab:ext-idioma} desglosa la calidad por idioma.
El \emph{silencio} es la fracción de los 12 hallazgos sin mención y no requiere etiquetas.
En v1, un reporte en español quedaba en silencio en el $73\%$ de los hallazgos, frente al $45\%$ en inglés.
La diferencia venía sobre todo del estilo de redacción, no del vocabulario: los reportes en español, griego y neerlandés suelen declarar lo normal en bloque (``ligamentos cruzados y colaterales dentro de límites normales'', ``meniscos conservados'').
v4 redujo el silencio en español de $0.717$ a $0.560$.
En turco el silencio subió ($0.617 \rightarrow 0.703$) porque se eliminaron negativos espurios de artrosis.
En v3, palabras de lado como \emph{iç}/\emph{dış} (interno/externo) junto a ``normal'' negaban la artrosis de ese compartimento aunque la frase hablara del menisco.
Ese aumento de silencio es, por tanto, una corrección.
El acuerdo por idioma en griego ($0.53$) y neerlandés ($0.63$) se apoya en solo 2--3 estudios gold, unas 30 decisiones, y no es estadísticamente concluyente.
Además, algunos reportes neerlandeses tienen la sección de hallazgos vacía; ese silencio proviene del dato y no del extractor.

\begin{table}[t]
  \caption{Silencio medio (fracción de hallazgos sin mención) y acuerdo con los gold, por idioma del reporte.}
  \label{tab:ext-idioma}
  \centering
  \small
  \begin{tabular}{@{}lrrcccccc@{}}
    \toprule
    & & & \multicolumn{4}{c}{Silencio} & \multicolumn{2}{c}{Acuerdo gold} \\
    \cmidrule(lr){4-7} \cmidrule(l){8-9}
    Idioma & Reportes & Gold & v1 & v2 & v3 & v4 & v1 & v4 \\
    \midrule
    Inglés     & 1736 & 28 & .445 & .393 & .393 & .383 & .827 & .812 \\
    Español    &  682 & 10 & .730 & .717 & .717 & .560 & .783 & .842 \\
    Turco      &  546 &  6 & .633 & .617 & .617 & .703 & .750 & .764 \\
    Croata     &  407 &  4 & .598 & .527 & .527 & .509 & .729 & .729 \\
    Griego     &  321 &  3 & .781 & .737 & .737 & .661 & .528 & .528 \\
    Alemán     &  262 &  2 & .589 & .543 & .543 & .513 & .792 & .792 \\
    Cirílico   &  220 &  3 & .717 & .678 & .678 & .649 & .750 & .750 \\
    Neerlandés &  151 &  2 & .726 & .670 & .670 & .604 & .542 & .625 \\
    Francés    &   81 &  0 & .384 & .340 & .340 & .294 & --   & --   \\
    \bottomrule
  \end{tabular}
\end{table}

\paragraph{Limitaciones.}
Los 58 gold cumplen dos funciones: miden la calidad del extractor y, más adelante, evalúan el modelo final.
Dos parámetros del extractor (derrame leve $=0.5$ y sinovitis leve $=1$) se fijaron mirando su efecto sobre esos gold, lo que puede hacer algo optimista la evaluación final.
Para acotar ese sesgo, el resto de las reglas son clínicas y generales: grados de Outerbridge, negación de estructuras en bloque y sinónimos.
Además, la selección de época y de umbrales del modelo se hace en un \emph{holdout} con etiquetas débiles, no en los gold.
   (requiere \usepackage{booktabs}, ya incluido en la plantilla)
% Fuente de los números: salida de MODE='labels' en Kaggle (4407 reportes, 58 estudios gold × 12 hallazgos = 696 decisiones).

\subsection{Extractor de etiquetas débiles: validación y evolución}
\label{sec:extractor}

Solo 58 de los 4407 estudios de entrenamiento tienen etiqueta experta (\emph{gold}).
El resto se etiqueta con un extractor propio de reglas que lee el reporte radiológico, escrito en al menos nueve idiomas.
Para cada hallazgo, el extractor busca menciones (término directo, o anatomía + patología a menos de 60 caracteres) y clasifica cada una según su contexto inmediato, al estilo de NegEx~\cite{chapman2001negex}:
negación $\rightarrow 0$, duda $\rightarrow 0.5$, edema óseo sin la palabra \emph{contusión} $\rightarrow 0.4$, condropatía de bajo grado $\rightarrow 0.4$, derrame leve $\rightarrow 0.5$, y afirmación $\rightarrow 1$.
El puntaje del estudio es el máximo de sus menciones.
Si no hay ninguna mención (\emph{silencio}), la etiqueta queda indefinida y en el entrenamiento recibe peso $0.3$.

\paragraph{Métricas.}
Cada versión se compara con los 58 gold.
El \emph{acuerdo} es la fracción de decisiones binarias (puntaje $\geq 0.5$, silencio $=0$) que coinciden con el experto, y la \emph{tasa de ruido} es $1-\text{acuerdo}$.
La \emph{mención} es la fracción de reportes gold en que el extractor encontró algo sobre el hallazgo.
El AUC usa el puntaje continuo y no depende del umbral, por eso es la métrica principal para comparar versiones.

\paragraph{Evolución.}
El extractor pasó por cuatro versiones (Tabla~\ref{tab:ext-cambios}).
Cada cambio salió de leer reportes reales con la función \texttt{show\_reports}, no de ajustar casos gold puntuales.
Cada regla nueva quedó fijada con una prueba de regresión: 27 en total.
El AUC macro subió de $0.739$ a $0.753$ y la tasa de ruido bajó de $0.226$ a $0.218$ (Tabla~\ref{tab:ext-versiones}).
La mención media pasó de $0.509$ a $0.566$ (Tabla~\ref{tab:ext-final}).

\begin{table}[t]
  \caption{Cambios del extractor por versión. Todos los cambios provienen de leer reportes reales.}
  \label{tab:ext-cambios}
  \centering
  \small
  \begin{tabular}{@{}lp{0.42\linewidth}p{0.42\linewidth}@{}}
    \toprule
    Versión & Cambio & Motivación \\
    \midrule
    v1 & Léxico multilingüe anatomía + patología; negación y duda estilo NegEx & Línea base \\
    v2 & Derrame y sinovitis leves $\rightarrow 0.3$; edema óseo $\rightarrow 0.4$; sinónimos de artrosis (cartílago, osteofitos, pinzamiento); sinovitis por engrosamiento sinovial y grasa de Hoffa; rotura meniscal por ``señal que alcanza la superficie''; puntaje del estudio = máximo
      & Especificidad de derrame $0.35$; sensibilidad de OA lateral y patelofemoral $<0.40$; mención de sinovitis $0.26$ \\
    v3 & Derrame leve $\rightarrow 0.5$ (incierto); sinovitis leve $\rightarrow 1$; condropatía grado I--II o superficial $\rightarrow 0.4$
      & En v2, los derrames leves resultaron positivos en el gold en proporción cercana a la mitad; grado I--II de Outerbridge no es artrosis \\
    v4 & Negación de estructuras nombradas en conjunto (``meniscos'', ``ligamentos colaterales'', ``cartílago'', ``médula ósea''); unión de líneas cortadas; ventana de negación de 35 a 60 caracteres; ``grado 2 y 3'' cuenta como alto grado; un lado anatómico solo (``menisco lateral íntegro'') ya no niega la artrosis de ese lado; ``osteocondral'' en griego deja de contar como artrosis
      & Reportes en español, griego y neerlandés que listan lo normal en bloque quedaban en silencio; falsos negativos y falsos positivos de artrosis \\
    \bottomrule
  \end{tabular}
\end{table}

\begin{table}[t]
  \caption{Acuerdo y AUC del extractor frente a los 58 gold, por versión. $n^{+}$: positivos según el experto.}
  \label{tab:ext-versiones}
  \centering
  \small
  \begin{tabular}{@{}lrcccccccc@{}}
    \toprule
    & & \multicolumn{4}{c}{Acuerdo} & \multicolumn{4}{c}{AUC} \\
    \cmidrule(lr){3-6} \cmidrule(l){7-10}
    Hallazgo & $n^{+}$ & v1 & v2 & v3 & v4 & v1 & v2 & v3 & v4 \\
    \midrule
    ACL              & 24 & .879 & .879 & .879 & .879 & .885 & .885 & .885 & .885 \\
    MCL              &  9 & .828 & .828 & .828 & .828 & .754 & .754 & .754 & .754 \\
    Medial Meniscus  & 26 & .776 & .776 & .776 & .759 & .768 & .768 & .768 & .752 \\
    Lateral Meniscus & 23 & .828 & .828 & .828 & .828 & .806 & .806 & .806 & .806 \\
    Medial OA        & 15 & .828 & .776 & .776 & .810 & .747 & .740 & .740 & .760 \\
    Lateral OA       & 11 & .810 & .793 & .810 & .810 & .635 & .725 & .733 & .733 \\
    PF OA            & 21 & .707 & .672 & .672 & .707 & .636 & .655 & .655 & .690 \\
    Effusion         & 35 & .638 & .621 & .638 & .638 & .666 & .666 & .666 & .666 \\
    Synovitis        & 27 & .655 & .621 & .672 & .672 & .636 & .655 & .658 & .658 \\
    Baker's          & 12 & .862 & .862 & .862 & .862 & .821 & .821 & .821 & .821 \\
    Contusion        & 19 & .690 & .793 & .793 & .793 & .751 & .760 & .760 & .760 \\
    Fracture         & 18 & .793 & .793 & .793 & .793 & .758 & .758 & .758 & .758 \\
    \midrule
    Media / macro    &    & .774 & .770 & .777 & .782 & .739 & .749 & .750 & .753 \\
    Tasa de ruido    &    & .226 & .230 & .223 & .218 &      &      &      &      \\
    \bottomrule
  \end{tabular}
\end{table}

\begin{table}[t]
  \caption{Versión final (v4) frente a los gold, con la mención de v1 como referencia.
    Sensibilidad y especificidad usan el umbral $0.5$.}
  \label{tab:ext-final}
  \centering
  \small
  \begin{tabular}{@{}lrcccccc@{}}
    \toprule
    & & \multicolumn{2}{c}{Mención} & & & & \\
    \cmidrule(lr){3-4}
    Hallazgo & $n^{+}$ & v1 & v4 & Acuerdo & Sens. & Espec. & AUC \\
    \midrule
    ACL              & 24 & .690 & .724 & .879 & .917 & .853 & .885 \\
    MCL              &  9 & .500 & .621 & .828 & .667 & .857 & .754 \\
    Medial Meniscus  & 26 & .517 & .621 & .759 & .692 & .812 & .752 \\
    Lateral Meniscus & 23 & .569 & .586 & .828 & .652 & .943 & .806 \\
    Medial OA        & 15 & .431 & .500 & .810 & .667 & .860 & .760 \\
    Lateral OA       & 11 & .448 & .448 & .810 & .636 & .851 & .733 \\
    PF OA            & 21 & .552 & .621 & .707 & .667 & .730 & .690 \\
    Effusion         & 35 & .845 & .845 & .638 & .829 & .348 & .666 \\
    Synovitis        & 27 & .259 & .362 & .672 & .444 & .871 & .658 \\
    Baker's          & 12 & .362 & .362 & .862 & .750 & .891 & .821 \\
    Contusion        & 19 & .517 & .603 & .793 & .684 & .846 & .760 \\
    Fracture         & 18 & .414 & .500 & .793 & .667 & .850 & .758 \\
    \midrule
    Media            &    & .509 & .566 & .782 & .689 & .809 & .753 \\
    \bottomrule
  \end{tabular}
\end{table}

\paragraph{Lectura por hallazgo.}
Los hallazgos que el reporte nombra de forma explícita y binaria son los mejor recuperados: ACL (AUC $0.885$), quiste de Baker ($0.821$) y menisco lateral ($0.806$).
El derrame muestra el límite de un umbral fijo.
De v1 a v3 su AUC no cambió ($0.666$), pero sensibilidad y especificidad se intercambiaron según el valor asignado al derrame \emph{leve}.
En v2 ese valor fue $0.3$: se perdieron 11 positivos del experto y se ganaron 10 negativos, es decir, los expertos consideran positiva aproximadamente la mitad de los derrames leves.
Por eso, desde v3 el derrame leve vale $0.5$ (incierto) y entra al entrenamiento con peso $0.5$.
La sinovitis tiene la mención más baja ($0.36$) aunque el experto la marca en 27 de 58 estudios.
Probablemente el experto la infiere de la imagen y no del texto, y ninguna regla sobre el reporte puede recuperar esa información.
v4 introdujo una regresión pequeña en el menisco medial: un estudio negativo pasó a positivo (especificidad de $0.844$ a $0.812$).

\paragraph{Sesgo lingüístico.}
La Tabla~\ref{tab:ext-idioma} desglosa la calidad por idioma.
El \emph{silencio} es la fracción de los 12 hallazgos sin mención y no requiere etiquetas.
En v1, un reporte en español quedaba en silencio en el $73\%$ de los hallazgos, frente al $45\%$ en inglés.
La diferencia venía sobre todo del estilo de redacción, no del vocabulario: los reportes en español, griego y neerlandés suelen declarar lo normal en bloque (``ligamentos cruzados y colaterales dentro de límites normales'', ``meniscos conservados'').
v4 redujo el silencio en español de $0.717$ a $0.560$.
En turco el silencio subió ($0.617 \rightarrow 0.703$) porque se eliminaron negativos espurios de artrosis.
En v3, palabras de lado como \emph{iç}/\emph{dış} (interno/externo) junto a ``normal'' negaban la artrosis de ese compartimento aunque la frase hablara del menisco.
Ese aumento de silencio es, por tanto, una corrección.
El acuerdo por idioma en griego ($0.53$) y neerlandés ($0.63$) se apoya en solo 2--3 estudios gold, unas 30 decisiones, y no es estadísticamente concluyente.
Además, algunos reportes neerlandeses tienen la sección de hallazgos vacía; ese silencio proviene del dato y no del extractor.

\begin{table}[t]
  \caption{Silencio medio (fracción de hallazgos sin mención) y acuerdo con los gold, por idioma del reporte.}
  \label{tab:ext-idioma}
  \centering
  \small
  \begin{tabular}{@{}lrrcccccc@{}}
    \toprule
    & & & \multicolumn{4}{c}{Silencio} & \multicolumn{2}{c}{Acuerdo gold} \\
    \cmidrule(lr){4-7} \cmidrule(l){8-9}
    Idioma & Reportes & Gold & v1 & v2 & v3 & v4 & v1 & v4 \\
    \midrule
    Inglés     & 1736 & 28 & .445 & .393 & .393 & .383 & .827 & .812 \\
    Español    &  682 & 10 & .730 & .717 & .717 & .560 & .783 & .842 \\
    Turco      &  546 &  6 & .633 & .617 & .617 & .703 & .750 & .764 \\
    Croata     &  407 &  4 & .598 & .527 & .527 & .509 & .729 & .729 \\
    Griego     &  321 &  3 & .781 & .737 & .737 & .661 & .528 & .528 \\
    Alemán     &  262 &  2 & .589 & .543 & .543 & .513 & .792 & .792 \\
    Cirílico   &  220 &  3 & .717 & .678 & .678 & .649 & .750 & .750 \\
    Neerlandés &  151 &  2 & .726 & .670 & .670 & .604 & .542 & .625 \\
    Francés    &   81 &  0 & .384 & .340 & .340 & .294 & --   & --   \\
    \bottomrule
  \end{tabular}
\end{table}

\paragraph{Limitaciones.}
Los 58 gold cumplen dos funciones: miden la calidad del extractor y, más adelante, evalúan el modelo final.
Dos parámetros del extractor (derrame leve $=0.5$ y sinovitis leve $=1$) se fijaron mirando su efecto sobre esos gold, lo que puede hacer algo optimista la evaluación final.
Para acotar ese sesgo, el resto de las reglas son clínicas y generales: grados de Outerbridge, negación de estructuras en bloque y sinónimos.
Además, la selección de época y de umbrales del modelo se hace en un \emph{holdout} con etiquetas débiles, no en los gold.
   (requiere \usepackage{booktabs}, ya incluido en la plantilla)
% Fuente de los números: salida de MODE='labels' en Kaggle (4407 reportes, 58 estudios gold × 12 hallazgos = 696 decisiones).

\subsection{Extractor de etiquetas débiles: validación y evolución}
\label{sec:extractor}

Solo 58 de los 4407 estudios de entrenamiento tienen etiqueta experta (\emph{gold}).
El resto se etiqueta con un extractor propio de reglas que lee el reporte radiológico, escrito en al menos nueve idiomas.
Para cada hallazgo, el extractor busca menciones (término directo, o anatomía + patología a menos de 60 caracteres) y clasifica cada una según su contexto inmediato, al estilo de NegEx~\cite{chapman2001negex}:
negación $\rightarrow 0$, duda $\rightarrow 0.5$, edema óseo sin la palabra \emph{contusión} $\rightarrow 0.4$, condropatía de bajo grado $\rightarrow 0.4$, derrame leve $\rightarrow 0.5$, y afirmación $\rightarrow 1$.
El puntaje del estudio es el máximo de sus menciones.
Si no hay ninguna mención (\emph{silencio}), la etiqueta queda indefinida y en el entrenamiento recibe peso $0.3$.

\paragraph{Métricas.}
Cada versión se compara con los 58 gold.
El \emph{acuerdo} es la fracción de decisiones binarias (puntaje $\geq 0.5$, silencio $=0$) que coinciden con el experto, y la \emph{tasa de ruido} es $1-\text{acuerdo}$.
La \emph{mención} es la fracción de reportes gold en que el extractor encontró algo sobre el hallazgo.
El AUC usa el puntaje continuo y no depende del umbral, por eso es la métrica principal para comparar versiones.

\paragraph{Evolución.}
El extractor pasó por cuatro versiones (Tabla~\ref{tab:ext-cambios}).
Cada cambio salió de leer reportes reales con la función \texttt{show\_reports}, no de ajustar casos gold puntuales.
Cada regla nueva quedó fijada con una prueba de regresión: 27 en total.
El AUC macro subió de $0.739$ a $0.753$ y la tasa de ruido bajó de $0.226$ a $0.218$ (Tabla~\ref{tab:ext-versiones}).
La mención media pasó de $0.509$ a $0.566$ (Tabla~\ref{tab:ext-final}).

\begin{table}[t]
  \caption{Cambios del extractor por versión. Todos los cambios provienen de leer reportes reales.}
  \label{tab:ext-cambios}
  \centering
  \small
  \begin{tabular}{@{}lp{0.42\linewidth}p{0.42\linewidth}@{}}
    \toprule
    Versión & Cambio & Motivación \\
    \midrule
    v1 & Léxico multilingüe anatomía + patología; negación y duda estilo NegEx & Línea base \\
    v2 & Derrame y sinovitis leves $\rightarrow 0.3$; edema óseo $\rightarrow 0.4$; sinónimos de artrosis (cartílago, osteofitos, pinzamiento); sinovitis por engrosamiento sinovial y grasa de Hoffa; rotura meniscal por ``señal que alcanza la superficie''; puntaje del estudio = máximo
      & Especificidad de derrame $0.35$; sensibilidad de OA lateral y patelofemoral $<0.40$; mención de sinovitis $0.26$ \\
    v3 & Derrame leve $\rightarrow 0.5$ (incierto); sinovitis leve $\rightarrow 1$; condropatía grado I--II o superficial $\rightarrow 0.4$
      & En v2, los derrames leves resultaron positivos en el gold en proporción cercana a la mitad; grado I--II de Outerbridge no es artrosis \\
    v4 & Negación de estructuras nombradas en conjunto (``meniscos'', ``ligamentos colaterales'', ``cartílago'', ``médula ósea''); unión de líneas cortadas; ventana de negación de 35 a 60 caracteres; ``grado 2 y 3'' cuenta como alto grado; un lado anatómico solo (``menisco lateral íntegro'') ya no niega la artrosis de ese lado; ``osteocondral'' en griego deja de contar como artrosis
      & Reportes en español, griego y neerlandés que listan lo normal en bloque quedaban en silencio; falsos negativos y falsos positivos de artrosis \\
    \bottomrule
  \end{tabular}
\end{table}

\begin{table}[t]
  \caption{Acuerdo y AUC del extractor frente a los 58 gold, por versión. $n^{+}$: positivos según el experto.}
  \label{tab:ext-versiones}
  \centering
  \small
  \begin{tabular}{@{}lrcccccccc@{}}
    \toprule
    & & \multicolumn{4}{c}{Acuerdo} & \multicolumn{4}{c}{AUC} \\
    \cmidrule(lr){3-6} \cmidrule(l){7-10}
    Hallazgo & $n^{+}$ & v1 & v2 & v3 & v4 & v1 & v2 & v3 & v4 \\
    \midrule
    ACL              & 24 & .879 & .879 & .879 & .879 & .885 & .885 & .885 & .885 \\
    MCL              &  9 & .828 & .828 & .828 & .828 & .754 & .754 & .754 & .754 \\
    Medial Meniscus  & 26 & .776 & .776 & .776 & .759 & .768 & .768 & .768 & .752 \\
    Lateral Meniscus & 23 & .828 & .828 & .828 & .828 & .806 & .806 & .806 & .806 \\
    Medial OA        & 15 & .828 & .776 & .776 & .810 & .747 & .740 & .740 & .760 \\
    Lateral OA       & 11 & .810 & .793 & .810 & .810 & .635 & .725 & .733 & .733 \\
    PF OA            & 21 & .707 & .672 & .672 & .707 & .636 & .655 & .655 & .690 \\
    Effusion         & 35 & .638 & .621 & .638 & .638 & .666 & .666 & .666 & .666 \\
    Synovitis        & 27 & .655 & .621 & .672 & .672 & .636 & .655 & .658 & .658 \\
    Baker's          & 12 & .862 & .862 & .862 & .862 & .821 & .821 & .821 & .821 \\
    Contusion        & 19 & .690 & .793 & .793 & .793 & .751 & .760 & .760 & .760 \\
    Fracture         & 18 & .793 & .793 & .793 & .793 & .758 & .758 & .758 & .758 \\
    \midrule
    Media / macro    &    & .774 & .770 & .777 & .782 & .739 & .749 & .750 & .753 \\
    Tasa de ruido    &    & .226 & .230 & .223 & .218 &      &      &      &      \\
    \bottomrule
  \end{tabular}
\end{table}

\begin{table}[t]
  \caption{Versión final (v4) frente a los gold, con la mención de v1 como referencia.
    Sensibilidad y especificidad usan el umbral $0.5$.}
  \label{tab:ext-final}
  \centering
  \small
  \begin{tabular}{@{}lrcccccc@{}}
    \toprule
    & & \multicolumn{2}{c}{Mención} & & & & \\
    \cmidrule(lr){3-4}
    Hallazgo & $n^{+}$ & v1 & v4 & Acuerdo & Sens. & Espec. & AUC \\
    \midrule
    ACL              & 24 & .690 & .724 & .879 & .917 & .853 & .885 \\
    MCL              &  9 & .500 & .621 & .828 & .667 & .857 & .754 \\
    Medial Meniscus  & 26 & .517 & .621 & .759 & .692 & .812 & .752 \\
    Lateral Meniscus & 23 & .569 & .586 & .828 & .652 & .943 & .806 \\
    Medial OA        & 15 & .431 & .500 & .810 & .667 & .860 & .760 \\
    Lateral OA       & 11 & .448 & .448 & .810 & .636 & .851 & .733 \\
    PF OA            & 21 & .552 & .621 & .707 & .667 & .730 & .690 \\
    Effusion         & 35 & .845 & .845 & .638 & .829 & .348 & .666 \\
    Synovitis        & 27 & .259 & .362 & .672 & .444 & .871 & .658 \\
    Baker's          & 12 & .362 & .362 & .862 & .750 & .891 & .821 \\
    Contusion        & 19 & .517 & .603 & .793 & .684 & .846 & .760 \\
    Fracture         & 18 & .414 & .500 & .793 & .667 & .850 & .758 \\
    \midrule
    Media            &    & .509 & .566 & .782 & .689 & .809 & .753 \\
    \bottomrule
  \end{tabular}
\end{table}

\paragraph{Lectura por hallazgo.}
Los hallazgos que el reporte nombra de forma explícita y binaria son los mejor recuperados: ACL (AUC $0.885$), quiste de Baker ($0.821$) y menisco lateral ($0.806$).
El derrame muestra el límite de un umbral fijo.
De v1 a v3 su AUC no cambió ($0.666$), pero sensibilidad y especificidad se intercambiaron según el valor asignado al derrame \emph{leve}.
En v2 ese valor fue $0.3$: se perdieron 11 positivos del experto y se ganaron 10 negativos, es decir, los expertos consideran positiva aproximadamente la mitad de los derrames leves.
Por eso, desde v3 el derrame leve vale $0.5$ (incierto) y entra al entrenamiento con peso $0.5$.
La sinovitis tiene la mención más baja ($0.36$) aunque el experto la marca en 27 de 58 estudios.
Probablemente el experto la infiere de la imagen y no del texto, y ninguna regla sobre el reporte puede recuperar esa información.
v4 introdujo una regresión pequeña en el menisco medial: un estudio negativo pasó a positivo (especificidad de $0.844$ a $0.812$).

\paragraph{Sesgo lingüístico.}
La Tabla~\ref{tab:ext-idioma} desglosa la calidad por idioma.
El \emph{silencio} es la fracción de los 12 hallazgos sin mención y no requiere etiquetas.
En v1, un reporte en español quedaba en silencio en el $73\%$ de los hallazgos, frente al $45\%$ en inglés.
La diferencia venía sobre todo del estilo de redacción, no del vocabulario: los reportes en español, griego y neerlandés suelen declarar lo normal en bloque (``ligamentos cruzados y colaterales dentro de límites normales'', ``meniscos conservados'').
v4 redujo el silencio en español de $0.717$ a $0.560$.
En turco el silencio subió ($0.617 \rightarrow 0.703$) porque se eliminaron negativos espurios de artrosis.
En v3, palabras de lado como \emph{iç}/\emph{dış} (interno/externo) junto a ``normal'' negaban la artrosis de ese compartimento aunque la frase hablara del menisco.
Ese aumento de silencio es, por tanto, una corrección.
El acuerdo por idioma en griego ($0.53$) y neerlandés ($0.63$) se apoya en solo 2--3 estudios gold, unas 30 decisiones, y no es estadísticamente concluyente.
Además, algunos reportes neerlandeses tienen la sección de hallazgos vacía; ese silencio proviene del dato y no del extractor.

\begin{table}[t]
  \caption{Silencio medio (fracción de hallazgos sin mención) y acuerdo con los gold, por idioma del reporte.}
  \label{tab:ext-idioma}
  \centering
  \small
  \begin{tabular}{@{}lrrcccccc@{}}
    \toprule
    & & & \multicolumn{4}{c}{Silencio} & \multicolumn{2}{c}{Acuerdo gold} \\
    \cmidrule(lr){4-7} \cmidrule(l){8-9}
    Idioma & Reportes & Gold & v1 & v2 & v3 & v4 & v1 & v4 \\
    \midrule
    Inglés     & 1736 & 28 & .445 & .393 & .393 & .383 & .827 & .812 \\
    Español    &  682 & 10 & .730 & .717 & .717 & .560 & .783 & .842 \\
    Turco      &  546 &  6 & .633 & .617 & .617 & .703 & .750 & .764 \\
    Croata     &  407 &  4 & .598 & .527 & .527 & .509 & .729 & .729 \\
    Griego     &  321 &  3 & .781 & .737 & .737 & .661 & .528 & .528 \\
    Alemán     &  262 &  2 & .589 & .543 & .543 & .513 & .792 & .792 \\
    Cirílico   &  220 &  3 & .717 & .678 & .678 & .649 & .750 & .750 \\
    Neerlandés &  151 &  2 & .726 & .670 & .670 & .604 & .542 & .625 \\
    Francés    &   81 &  0 & .384 & .340 & .340 & .294 & --   & --   \\
    \bottomrule
  \end{tabular}
\end{table}

\paragraph{Limitaciones.}
Los 58 gold cumplen dos funciones: miden la calidad del extractor y, más adelante, evalúan el modelo final.
Dos parámetros del extractor (derrame leve $=0.5$ y sinovitis leve $=1$) se fijaron mirando su efecto sobre esos gold, lo que puede hacer algo optimista la evaluación final.
Para acotar ese sesgo, el resto de las reglas son clínicas y generales: grados de Outerbridge, negación de estructuras en bloque y sinónimos.
Además, la selección de época y de umbrales del modelo se hace en un \emph{holdout} con etiquetas débiles, no en los gold.
